\clearpage
\onecolumn

\section*{Appendix}

\subsection*{Prism-Effect Category Synthesis}

Table~\ref{tab:article-appendix-effect-categories} expands the perturbation categories introduced in the related work.

\begin{table}[htbp]
    \centering
    \scriptsize
    \setlength{\tabcolsep}{4pt}
    \renewcommand{\arraystretch}{1.15}
    \begin{tabular}{|
        p{0.16\textwidth}|
        p{0.29\textwidth}|
        p{0.48\textwidth}|}
        \hline
        \textbf{Icon} & \textbf{Name} & \textbf{Description} \\
        \hline
        \begin{minipage}[t]{\linewidth}
            \vspace{0pt}
            \centering
            \includegraphics[width=2.7cm]{img/1_Background/Icons/Skewness.png}
        \end{minipage}
        & \begin{minipage}[t]{\linewidth}
            \vspace{0pt}
            \textbf{Shifted rendering or scene}
        \end{minipage}
        & \begin{minipage}[t]{\linewidth}
            \vspace{0pt}
            A manipulation where the visual scene, rendering geometry, or visible response representation is displaced relative to the participant. This is conceptually closest to conventional optical prism displacement when the whole visual field is shifted, although VR studies differ in whether they shift the scene, the camera frustum, or a task-specific visual reference \cite{Ramos2019,anan2025comparative}.
        \end{minipage} \\
        \hline
        \begin{minipage}[t]{\linewidth}
            \vspace{0pt}
            \centering
            \includegraphics[width=2.7cm]{img/1_Background/Icons/Translate - Hand.png}
        \end{minipage}
        & \begin{minipage}[t]{\linewidth}
            \vspace{0pt}
            \textbf{Virtual effector or response displacement}
        \end{minipage}
        & \begin{minipage}[t]{\linewidth}
            \vspace{0pt}
            A mismatch where the rendered hand, controller, endpoint, or response origin is offset relative to the tracked real movement. Several VR prism-adaptation prototypes use this effector-displacement logic rather than physically shifting the entire visual field \cite{kim2017development,cho2022feasibility,Bourgeois2021}.
        \end{minipage} \\
        \cline{1-2}
        \begin{minipage}[t]{\linewidth}
            \vspace{0pt}
            \centering
            \includegraphics[width=2.7cm]{img/1_Background/Icons/Translate - Arm.png}
        \end{minipage}
        & \begin{minipage}[t]{\linewidth}
            \vspace{0pt}
            with Arm
        \end{minipage}
        & \\
        \cline{1-2}
        \begin{minipage}[t]{\linewidth}
            \vspace{0pt}
            \centering
            \includegraphics[width=2.7cm]{img/1_Background/Icons/Translate - Controller.png}
        \end{minipage}
        & \begin{minipage}[t]{\linewidth}
            \vspace{0pt}
            with Controller
        \end{minipage}
        & \\
        \hline
        \begin{minipage}[t]{\linewidth}
            \vspace{0pt}
            \centering
            \includegraphics[width=2.7cm]{img/1_Background/Icons/Rotation - Controller v.1.png}
        \end{minipage}
        & \begin{minipage}[t]{\linewidth}
            \vspace{0pt}
            \textbf{Visuomotor rotation}
        \end{minipage}
        & \begin{minipage}[t]{\linewidth}
            \vspace{0pt}
            A manipulation where the direction of visual feedback is rotated relative to the executed movement. Rotation is a standard visuomotor-adaptation manipulation and is useful as a comparison case because it changes the action-feedback mapping differently from a constant lateral offset \cite{gerken2025sense}.
        \end{minipage} \\
        \cline{1-2}
        \begin{minipage}[t]{\linewidth}
            \vspace{0pt}
            \centering
            \includegraphics[width=2.7cm]{img/1_Background/Icons/Rotation - Hand.png}
        \end{minipage}
        & \begin{minipage}[t]{\linewidth}
            \vspace{0pt}
            with Hand
        \end{minipage}
        & \\
        \cline{1-2}
        \begin{minipage}[t]{\linewidth}
            \vspace{0pt}
            \centering
            \includegraphics[width=2.7cm]{img/1_Background/Icons/Rotation - Arm.png}
        \end{minipage}
        & \begin{minipage}[t]{\linewidth}
            \vspace{0pt}
            with Arm
        \end{minipage}
        & \\
        \hline
    \end{tabular}
    \caption{Conceptual prism-effect categories used to organize the reviewed VR perturbations. Blue elements represent virtual components, grey elements represent physical components, dotted lines denote conceptual response or raycast directions, and red circles indicate target locations.}
    \label{tab:article-appendix-effect-categories}
\end{table}

\clearpage
\subsection*{Measurement Task Category Synthesis}

Table~\ref{tab:article-appendix-task-categories} expands the prism-therapy measurement task categories introduced in the related work.

\begin{table}[htbp]
    \centering
    \footnotesize
    \setlength{\tabcolsep}{4pt}
    \renewcommand{\arraystretch}{1.15}
    \begin{tabular}{|
        p{0.18\textwidth}|
        p{0.24\textwidth}|
        p{0.51\textwidth}|}
        \hline
        \textbf{Icon} & \textbf{Name} & \textbf{Description} \\
        \hline
        \hfill\includegraphics[width=2.5cm]{img/1_Background/Icons/Open Loop Pointing.png}\hfill\null
        & Open-loop pointing
        & A post-exposure pointing measure used to estimate the direction and magnitude of the sensorimotor after-effect after the perturbation has been removed \cite{Rossetti1998Prism,Ramos2019,anan2025comparative}. \\
        \hline
        \hfill\includegraphics[width=2.5cm]{img/1_Background/Icons/Line Bisection.png}\hfill\null
        & Line bisection
        & A spatial judgement task in which the participant indicates the perceived midpoint of a horizontal line. Baseline-to-post changes can indicate transfer from the trained exposure movement to a spatial estimation task \cite{frassinetti2002long,Serino2011,Bourgeois2021}. \\
        \hline
        \hfill\includegraphics[width=2.5cm]{img/1_Background/Icons/Landmark Bisection.png}\hfill\null
        & Landmark judgement
        & A perceptual judgement task in which the participant decides whether a pre-bisected line is transected to the left or right of the true midpoint. Because it does not require the same pointing endpoint measure, it can help separate perceptual bias from motor execution \cite{Serino2011,Bourgeois2021,anan2025comparative}. \\
        \hline
    \end{tabular}
    \caption{Task categories used to describe behavioural assessment in the reviewed literature and in the prototype design.}
    \label{tab:article-appendix-task-categories}
\end{table}

\clearpage
\subsection*{Exposure Quality Examples}

Figure~\ref{fig:article-appendix-quality} shows the exposure-quality examples referenced in the results. The two panels are kept in the appendix so the Shiny dashboard screenshots remain readable at full width.

\begin{figure}[htbp]
    \centering
    \includegraphics[width=0.95\linewidth]{img/5_Results/Quality1.png}
    \textbf{(a)} Low-hit-rate exposure

    \vspace{0.8em}

    \includegraphics[width=0.95\linewidth]{img/5_Results/Quality2.png}
    \textbf{(b)} Higher-hit-rate exposure
    \caption{Exposure quality examples. Both modules contain 60 attempts, but the hit rate and distance-to-target profiles differ substantially. The low-hit-rate block contains more misses and larger target-relative distances, while the higher-hit-rate block stays closer to the target after one large early miss.}
    \label{fig:article-appendix-quality}
\end{figure}

\clearpage
\twocolumn
